\clearpage
\appendix

\section{中心胞元表示定理的证明}
\label{app:center-cell}

\begin{proof}[定理~\ref{thm:center-cell-representation} 的证明]
  \emph{Rayleigh 系数。}由于 $\Omega_0$ 的闭包含于 $\mathcal C_0$ 内部，在 $\Gamma_-$ 与
  $\Omega_0$ 之间、$\Omega_0$ 与 $\Gamma_+$ 之间各有一条背景竖直带，$u$ 在其中按
  Rayleigh 波展开。把 $u$ 与 $\partial_xu$ 在 $\Gamma_\pm$ 上的迹投影到 $\psi_m$，每个
  模态的两个振幅由矩阵 $\bigl[\begin{smallmatrix}1&1\\ \ii\gamma_m&-\ii\gamma_m
  \end{smallmatrix}\bigr]$ 唯一确定（$\gamma_m\ne0$）。取左带中的右行振幅为
  $a^-_{0,m}$（锚定在 $x_-$）、右带中的左行振幅为 $b^+_{0,m}$（锚定在 $x_+$），由此得到
  $u_{\mathrm{hom}}$，系数列唯一。\pending{2026-11-06 审稿意见 I6：补级数收敛的论证}
  记 $f^\pm_m,g^\pm_m$ 为 $u|_{\Gamma_\pm}$ 与 $\partial_xu|_{\Gamma_\pm}$ 关于 $\psi_m$ 的
  系数，则
  $a^-_{0,m}=\tfrac12\bigl(f^-_m+g^-_m/(\ii\gamma_m)\bigr)$，
  $b^+_{0,m}=\tfrac12\bigl(f^+_m-g^+_m/(\ii\gamma_m)\bigr)$。记
  $\langle\beta_m\rangle:=(1+\beta_m^2)^{1/2}$。$u$ 在端口附近属于 $H^1$，故
  $\sum_m\langle\beta_m\rangle|f^\pm_m|^2<\infty$、
  $\sum_m\langle\beta_m\rangle^{-1}|g^\pm_m|^2<\infty$；由假设~\ref{ass:no-wood} 及
  $|\gamma_m|\sim|\beta_m|$（$|m|\to\infty$），存在 $c>0$ 使 $|\gamma_m|\geq c\langle\beta_m\rangle$，
  于是 $\sum_m\langle\beta_m\rangle(|a^-_{0,m}|^2+|b^+_{0,m}|^2)<\infty$。另一方面，
  $r_m^{x_-}$ 在 $x\geq x_-$、$l_m^{x_+}$ 在 $x\leq x_+$ 上不增长，且
  $|\gamma_m|\leq\langle\beta_m\rangle+\kappa_{\mathrm e}$；对凋落模态
  $\operatorname{Im}\gamma_m\geq c'\langle\beta_m\rangle$，故
  $\lVert r_m^{x_-}\rVert_{H^1(\mathcal C_0)}^2,\lVert l_m^{x_+}\rVert_{H^1(\mathcal C_0)}^2
  \leq C\langle\beta_m\rangle$。不同 $m$ 的项及其一阶导数在每条竖直线上关于 $y$ 正交，所以
  式~\eqref{eq:center-homogeneous-field} 的两个级数在 $H^1(\mathcal C_0)$ 中收敛。若 $\Omega_0=\varnothing$，整个 $\mathcal C_0$ 是一条
  背景带，$u=u_{\mathrm{hom}}$，结论成立。以下设 $\Omega_0\ne\varnothing$。

  \emph{辅助问题。}在 $\partial\Omega_0$ 上，上标 $+$、$-$ 分别表示从柱外、柱内取迹，
  $\nu$ 由内指向外。令
  \begin{align*}
    w^+&:=u-u_{\mathrm{hom}}\quad\text{于 }\mathcal C_0\setminus\overline{\Omega_0},
    &w^-&:=u\quad\text{于 }\Omega_0.
  \end{align*}
  按构造，$w^+$ 在 $\Gamma_-$ 处只有左行分量、在 $\Gamma_+$ 处只有右行分量，即在两个
  端口都是出射的。考虑交换内外波数的全平面问题：在 $\Omega_0$ 内用
  $\kappa_{\mathrm e}$，在 $\RR^2\setminus\overline{\Omega_0}$ 内用 $\kappa_{\mathrm c}$
  并满足 Sommerfeld 出射条件，界面跳跃取为
  \begin{align*}
    v^+-v^-&=-w^+-w^-,
    &\partial_\nu v^+-\partial_\nu v^-&=-\partial_\nu w^+-\partial_\nu w^-.
  \end{align*}
  由假设它有唯一解 $v$。定义
  \begin{align*}
    \tau&:=w^+-v^-=-w^--v^+,
    &\omega&:=-\partial_\nu w^++\partial_\nu v^-=\partial_\nu w^-+\partial_\nu v^+,
  \end{align*}
  两种写法的相等分别来自上面两个跳跃条件。

  \emph{内部表示。}对 $\mathbf z\in\Omega_0$，$w^-$ 的内部 Green 表示给出
  $\mathcal S^{(\kappa_{\mathrm c})}[\partial_\nu w^-]-\mathcal D^{(\kappa_{\mathrm c})}[w^-]=w^-$；
  $v^+$ 是 $\kappa_{\mathrm c}$ 的出射外部解，其外部 Green 公式在柱内给出
  $\mathcal S^{(\kappa_{\mathrm c})}[\partial_\nu v^+]-\mathcal D^{(\kappa_{\mathrm c})}[v^+]=0$。
  代入 $\tau,\omega$ 的第二种写法并相加，得
  $\mathcal S^{(\kappa_{\mathrm c})}[\omega]+\mathcal D^{(\kappa_{\mathrm c})}[\tau]=w^-$。

  \emph{外部表示。}在 $\mathcal C_0\setminus\overline{\Omega_0}$ 上对 $w^+$ 与
  $G^{(\kappa_{\mathrm e})}_\QP(\mathbf z,\cdot)$ 用 Green 第二恒等式：顶、底边界的贡献
  因准周期相位相消；在两个端口上，$w^+$ 与 $G_\QP(\mathbf z,\cdot)$ 按模态都是同向出射
  的，逐模 Wronskian 为零，贡献也相消。因此
  $-\mathcal S^{(\kappa_{\mathrm e})}_\QP[\partial_\nu w^+]
  +\mathcal D^{(\kappa_{\mathrm e})}_\QP[w^+]=w^+$。另一方面，$v^-$ 在 $\Omega_0$ 内是
  $\kappa_{\mathrm e}$ 的解，而对柱外的 $\mathbf z$，$G_\QP(\mathbf z,\cdot)$ 在
  $\Omega_0$ 内光滑（$\mathbf z$ 的 $y$ 向像点都不在 $\Omega_0$ 中），故
  $\mathcal S^{(\kappa_{\mathrm e})}_\QP[\partial_\nu v^-]
  -\mathcal D^{(\kappa_{\mathrm e})}_\QP[v^-]=0$。代入 $\tau,\omega$ 的第一种写法并相加，
  得 $\mathcal S^{(\kappa_{\mathrm e})}_\QP[\omega]+\mathcal D^{(\kappa_{\mathrm e})}_\QP[\tau]
  =w^+$。加回 $u_{\mathrm{hom}}$ 即得外部表示。

  \emph{唯一性。}\pending{2026-11-06 审稿意见 I6：改为由交换问题的唯一性推出}若两组表示
  给出同一个场，则其差表示零场。由第一段的端口系数提取，Rayleigh 系数之差为零。记密度之差
  仍为 $\tau,\omega$，$L_{\mathrm e}:=\mathcal S^{(\kappa_{\mathrm e})}_\QP[\omega]
  +\mathcal D^{(\kappa_{\mathrm e})}_\QP[\tau]$，
  $L_{\mathrm c}:=\mathcal S^{(\kappa_{\mathrm c})}[\omega]+\mathcal D^{(\kappa_{\mathrm c})}[\tau]$，
  则 $L_{\mathrm e}$ 在柱外、$L_{\mathrm c}$ 在柱内为零。由跳跃关系（$G_\QP-G$ 在
  $\partial\Omega_0$ 附近光滑，跳跃与自由空间层势相同），
  $L_{\mathrm e}^-=-\tau$、$\partial_\nu L_{\mathrm e}^-=\omega$，
  $L_{\mathrm c}^+=\tau$、$\partial_\nu L_{\mathrm c}^+=-\omega$。令 $v^-:=L_{\mathrm e}$
  于 $\Omega_0$、$v^+:=-L_{\mathrm c}$ 于 $\RR^2\setminus\overline{\Omega_0}$，则 $v$ 是
  交换问题的解，出射，且两种跳跃都为零。由交换问题的唯一性 $v=0$，故
  $\tau=-v^-=0$、$\omega=\partial_\nu v^-=0$。这是 Barnett 与 Greengard 的互补场论证
  \cite[附录 A]{BarnettGreengard2010}（源于 \cite{ColtonKress1983}）在中心胞元上的版本；
  这里多出的是端口 Rayleigh 项，它已在上一步由系数提取消去。

  \emph{交换问题的唯一性。}设 $v$ 满足交换问题且跳跃为零。在 $\Omega_0$ 上用 Green 第一
  恒等式，
  \begin{align*}
    \int_{\partial\Omega_0}\overline{v^-}\,\partial_\nu v^-\dd s
    =\int_{\Omega_0}(|\nabla v|^2-\kappa_{\mathrm e}^2|v|^2)\dd\mathbf z
  \end{align*}
  为实数。迹相等，故
  $\int_{\partial\Omega_0}\overline{v^+}\,\partial_\nu v^+\dd s$ 也为实数。在
  $B_R\setminus\overline{\Omega_0}$ 上用 Green 第一恒等式，得
  \begin{align*}
    \operatorname{Im}\int_{\partial B_R}\overline v\,\partial_rv\dd s=0 .
  \end{align*}
  由 Rellich 引理
  \cite{ColtonKress1983}，$v^+=0$。于是 $v^-$ 的 Cauchy 数据为零，把它在柱外延拓为零得到
  $\kappa_{\mathrm e}$ 的 Helmholtz 解，由唯一延拓 $v^-=0$。
\end{proof}

\pending{2026-11-06 审稿意见 M3：第 (2) 条的密度空间与区域写精确}\pending{2026-11-06 新增：近像分解推论，正文第~\ref{subsec:layer-potentials}~小节引用}
下面的推论说明，定理中的准周期层势可以改写为近像层势加光滑修正，这正是实现中评价场的结构
（式~\eqref{eq:mfs-field}）。记号同第~\ref{subsec:layer-potentials}~小节。

\begin{corollary}[近像分解]
  \label{cor:near-image}
  设假设~\ref{ass:no-wood} 成立，$\kappa=\kappa_{\mathrm e}$。
  \begin{enumerate}
    \item 式~\eqref{eq:near-image-split} 定义的 $G_{\mathrm{far}}$ 在
      $E:=\RR^2\setminus\{nd\,e_y:|n|\geq2\}$ 上实解析，并满足
      $(\Delta+\kappa^2)G_{\mathrm{far}}=0$。
    \item 设 $\overline{\Omega}\subset\RR\times I_d$，$\tau\in H^{1/2}(\partial\Omega)$，
      $\omega\in H^{-1/2}(\partial\Omega)$，$\delta:=d/2-\max_{\partial\Omega}|y|>0$，
      $S_\delta:=\RR\times(-d/2-\delta,d/2+\delta)$。则在
      $S_\delta\setminus\bigl(\partial\Omega\cup(\partial\Omega\pm d\,e_y)\bigr)$ 上有
      \begin{align*}
        \mathcal S_\QP[\omega]+\mathcal D_\QP[\tau]
        &=\sum_{n=-1}^{1}q^n\bigl(\mathcal S[\omega]+\mathcal D[\tau]\bigr)(\cdot-nd\,e_y)
          +u_{\mathrm{corr}},
      \end{align*}
      其中 $u_{\mathrm{corr}}(\mathbf z):=\int_{\partial\Omega}\bigl(
      G_{\mathrm{far}}(\mathbf z-\mathbf z')\omega(\mathbf z')
      +\partial_{\nu'}[G_{\mathrm{far}}(\mathbf z-\mathbf z')]\tau(\mathbf z')\bigr)\dd s'$
      （积分理解为对偶配对）在整个 $S_\delta$ 上光滑并满足 Helmholtz 方程。
    \item 特别地，在定理~\ref{thm:center-cell-representation} 的条件下，$u$ 在
      $\mathcal C_0\setminus\overline{\Omega_0}$ 中等于近像层势、$u_{\mathrm{corr}}$ 与
      $u_{\mathrm{hom}}$ 三项之和，其中 $u_{\mathrm{corr}}$ 可延拓为闭胞元某邻域内的
      Helmholtz 解。
  \end{enumerate}
\end{corollary}

\begin{proof}
  (1) $G_\QP(\mathbf z,\mathbf z')$ 只依赖于 $\mathbf z-\mathbf z'$，且在分布意义下
  $(\Delta+\kappa^2)G_\QP(\cdot,\mathbf 0)=-\sum_{n\in\ZZ}q^n\delta_{nd\,e_y}$
  \cite{Linton1998}；而 $(\Delta+\kappa^2)G(\cdot,nd\,e_y)=-\delta_{nd\,e_y}$。相减后
  $G_{\mathrm{far}}$ 在 $E$ 上满足齐次 Helmholtz 方程，由常系数椭圆方程的内部正则性，它在
  $E$ 上实解析。

  (2) 把式~\eqref{eq:near-image-split} 代入层势的定义，并注意
  $\int_{\partial\Omega}G(\mathbf z,\mathbf z'+nd\,e_y)\omega(\mathbf z')\dd s'
  =\mathcal S[\omega](\mathbf z-nd\,e_y)$（双层势同理），即得分解式。若 $\mathbf z$ 在所给
  带内、$\mathbf z'\in\partial\Omega$，则 $|y-y'|<(d/2+\delta)+(d/2-\delta)=d$，故 $\mathbf z-\mathbf z'$ 与
  $\{nd\,e_y:|n|\geq2\}$ 的距离大于 $d$。由 (1)，核 $G_{\mathrm{far}}(\mathbf z-\mathbf z')$ 在
  $S_\delta\times\partial\Omega$ 的邻域上联合光滑，且对 $\mathbf z$ 满足 Helmholtz 方程；
  它与 $H^{\pm1/2}$ 中密度的对偶配对可以在配对下对 $\mathbf z$ 求导，故
  $u_{\mathrm{corr}}$ 在 $S_\delta$ 上光滑并满足 Helmholtz 方程。

  (3) 由定理~\ref{thm:center-cell-representation}，在 $\mathcal C_0\setminus\overline{\Omega_0}$
  中 $u=\mathcal S_\QP[\omega]+\mathcal D_\QP[\tau]+u_{\mathrm{hom}}$，对层势部分用 (2)。
\end{proof}

由推论，精确场在背景区域中与评价场~\eqref{eq:mfs-field} 具有相同的结构；Rayleigh 背景单独保留，
代理点源只需逼近 $u_{\mathrm{corr}}$ 这类在闭胞元邻域内光滑的解。本文不依赖这种逼近的收敛
理论：评价场的全部误差都由它的界面跳跃计量。

